\documentclass[11pt]{article}
\usepackage{amsmath, amssymb, amsthm}
\usepackage{geometry}
\usepackage{hyperref}
\usepackage{booktabs}
\usepackage{xcolor}
\usepackage{listings}

\geometry{margin=1.1in}

\newtheorem{lemma}{Lemma}

\newcommand{\eps}{\varepsilon}
\newcommand{\dt}{\Delta t}
\newcommand{\dx}{\Delta x}
\newcommand{\h}{h}
\newcommand{\code}[1]{\texttt{#1}}

\lstset{
  basicstyle=\ttfamily\small,
  columns=fullflexible,
  keepspaces=true,
  breaklines=true,
  frame=single,
}

\title{Non-Uniform Time Grid: Refining the First and Last Cell}
\author{Optimal Transport / Schr\"odinger Bridge Project}
\date{}

\begin{document}
\maketitle

\begin{abstract}
This note specifies the simplest non-trivial non-uniform time grid we consider for the
1D Python solver: take the uniform $n_t$-cell grid on $[0,1]$ and split only the first
and last cell in half. It sets out the two grids the LADMM splitting actually uses --
a staggered grid carrying the hard Fokker--Planck (FP) constraint, and a collocated grid
carrying the kinetic-energy (KE) prox -- and derives the full discrete FP constraint on
the refined-ends grid, row by row through both junctions. The surrounding solver
(\code{operators\_nu.py}, \code{projection\_nu.py}, \code{prox\_nu.py},
\code{pipeline\_nu.py}) is already generic over any strictly increasing edge sequence, so
this grid requires no new solver code, only a new grid constructor -- implemented as
\code{refine\_ends\_time\_grid} in \code{time\_grid.py} and regression-checked in
\code{scripts/validate\_nu.py} (5/5 checks pass, including this grid).
\end{abstract}

%------------------------------------------------------------------------------
\section{Two grids in the LADMM split}
\label{sec:two-grids}
%------------------------------------------------------------------------------

The LADMM splitting (\code{pipeline\_nu.py}) carries two copies of the state, on two
different grids, tied together by a consensus map:

\begin{itemize}
  \item \textbf{Constraint grid} (code: the \emph{x}-variable $(\mu,\psi)$). This is
        where the FP constraint is enforced \emph{exactly}, via the projection
        \code{proj\_fokker\_planck\_banded\_nu}. Staggered: density at interior time
        edges $\times$ space cell centers,
        \[
          \mu(t_j,\xi_k), \qquad j=1,\ldots,n_t^{\mathrm{ref}}-1,\quad k=1,\ldots,n_x,
        \]
        and momentum at time cell centers $\times$ interior space nodes,
        \[
          \psi(\tau_i,x_k), \qquad i=0,\ldots,n_t^{\mathrm{ref}}-1,\quad k=1,\ldots,n_x-1,
        \]
        with $\mu(t_0,\cdot)=\rho_0$, $\mu(t_{n_t^{\mathrm{ref}}},\cdot)=\rho_1$ (fixed
        marginals) and $\psi(\cdot,x_0)=\psi(\cdot,x_{n_x})=0$ (no-flux).
  \item \textbf{KE grid} (code: the \emph{y}-variable $(\rho,m)$). This is where the
        kinetic-energy prox \code{prox\_ke\_cc\_nu} acts, pointwise in $m^2/(2\rho)$, so
        $\rho$ and $m$ must be \emph{collocated} -- both at cell centers, both axes:
        \[
          \rho(\tau_i,\xi_k),\ m(\tau_i,\xi_k), \qquad i=0,\ldots,n_t^{\mathrm{ref}}-1,\quad
          k=1,\ldots,n_x.
        \]
\end{itemize}
Space is uniform throughout ($n_x$ cells, $\dx=1/n_x$, edges $x_k=k\dx$, centers
$\xi_k=(k-\tfrac12)\dx$) -- only the time axis is refined here, and only at the ends.

\subsection*{The consensus map}

\code{A\_fn} moves $(\mu,\psi)$ onto the KE grid's locations:
\begin{equation}
  \label{eq:consensus}
  \hat\rho(\tau_i,\xi_k) = \tfrac12\bigl[\mu(t_i,\xi_k)+\mu(t_{i+1},\xi_k)\bigr],
  \qquad
  \hat m(\tau_i,\xi_k) = \tfrac12\bigl[\psi(\tau_i,x_{k-1})+\psi(\tau_i,x_k)\bigr],
\end{equation}
and the LADMM iteration drives $\hat\rho\to\rho$, $\hat m\to m$ while $\mu,\psi$ stays
exactly FP-feasible and $\rho,m$ stays exactly KE-optimal. Note $\hat\rho$ only
interpolates \emph{in time} (both sides share the same $\xi_k$) and $\hat m$ only
interpolates \emph{in space} (both sides share the same $\tau_i$) -- $\psi$'s time
location is $\tau_i$ on \emph{both} grids already, so refining the time grid never
touches the momentum's time placement, only the density's.

\begin{lemma}[Midpoint interpolation is exact regardless of cell width]
\label{lem:midpoint}
Let $t_i<t_{i+1}$ be any two grid edges and $\tau_i:=\tfrac12(t_i+t_{i+1})$ their
midpoint (as \emph{every} cell center in this codebase is, by construction --
\code{time\_grid.\_time\_grid\_from\_edges}). For any values $v_i,v_{i+1}$, let
$L(t)$ be the unique affine function with $L(t_i)=v_i$, $L(t_{i+1})=v_{i+1}$. Then
$L(\tau_i)=\tfrac12(v_i+v_{i+1})$, independent of $\dt_i=t_{i+1}-t_i$.
\end{lemma}
\begin{proof}
$L(t)=v_i+(v_{i+1}-v_i)\frac{t-t_i}{\dt_i}$. At $t=\tau_i$, $t-t_i=\dt_i/2$, so
$L(\tau_i)=v_i+\tfrac12(v_{i+1}-v_i)=\tfrac12(v_i+v_{i+1})$.
\end{proof}
So the plain average in \eqref{eq:consensus}'s $\hat\rho$ formula is exactly the value
of the linear interpolant of $\mu$ \emph{at the cell center}, for \emph{any} $\dt_i$ --
this is the reason \code{interp\_t\_at\_phi} needs no change on a non-uniform grid at
all (\code{operators\_nu.py} reuses it verbatim from \code{operators.py}), and it is
the key fact behind the derivation in Section~\ref{sec:fp}.

%------------------------------------------------------------------------------
\section{The refined-ends grid}
%------------------------------------------------------------------------------

Take the uniform $n_t$-cell grid and split \emph{only} cell $0$ ($[t_0,t_1]$) and cell
$n_t-1$ ($[t_{n_t-1}, t_{n_t}]$) in half. Everything in between is untouched:
\[
\begin{array}{c|ccccccccc}
\text{edges} & 0 & \h/2 & \h & 2\h & 3\h & \cdots & (n_t{-}1)\h & (n_t{-}1)\h + \h/2 & 1 \\
\hline
\text{cell width} & \multicolumn{2}{c}{\h/2} & \multicolumn{5}{c}{\h} & \multicolumn{2}{c}{\h/2}
\end{array}
\]
So, writing $\h = 1/n_t$ for the \emph{base} resolution being refined,
\begin{equation}
  \label{eq:edges}
  \mathbf{t} = \bigl[\,0,\ \tfrac{\h}{2},\ \h,\ 2\h,\ 3\h,\ \ldots,\ (n_t{-}1)\h,\
  (n_t{-}1)\h + \tfrac{\h}{2},\ 1\,\bigr],
\end{equation}
i.e.\ one extra edge inserted into the first cell and one into the last. This has
$n_t+3$ edges, hence
\begin{equation}
  n_t^{\mathrm{ref}} = n_t + 2 \quad \text{cells,} \qquad
  \Delta t_i =
  \begin{cases}
    \h/2, & i \in \{0,1,\,n_t^{\mathrm{ref}}{-}2,\, n_t^{\mathrm{ref}}{-}1\}, \\
    \h,   & \text{otherwise,}
  \end{cases}
\end{equation}
and $\sum_i \Delta t_i = 2\cdot\tfrac{\h}{2} + (n_t-2)\h + 2\cdot\tfrac{\h}{2} = n_t \h = 1$,
as required. Cell centers and $\dt_{\mathrm{ref}} := 1/n_t^{\mathrm{ref}}$ follow the usual
construction from $\mathbf{t}$ (Eq.~\ref{eq:edges}); nothing else about the grid changes.
This affects \emph{both} grids of Section~\ref{sec:two-grids} identically: the KE grid's
$\tau_i$ are exactly the constraint grid's cell centers, so both move together, and only
$\mu$ (not $\psi$) picks up any new time-interpolation behavior, per the remark above.

\subsection{Worked example: $n_t = 6$}
\label{sec:worked}

Here $\h = 1/6$, $n_t^{\mathrm{ref}} = 8$.

\begin{center}
\begin{tabular}{c|cccccccc}
\toprule
$i$ (cell) & 0 & 1 & 2 & 3 & 4 & 5 & 6 & 7 \\
\midrule
$\Delta t_i$ & $1/12$ & $1/12$ & $1/6$ & $1/6$ & $1/6$ & $1/6$ & $1/12$ & $1/12$ \\
$\tau_i$ (center) & $1/24$ & $1/8$ & $1/4$ & $5/12$ & $7/12$ & $3/4$ & $7/8$ & $23/24$ \\
\bottomrule
\end{tabular}
\end{center}

\begin{center}
\begin{tabular}{c|ccccccccc}
\toprule
$j$ (edge) & 0 & 1 & 2 & 3 & 4 & 5 & 6 & 7 & 8 \\
\midrule
$t_j$ & $0$ & $1/12$ & $1/6$ & $1/3$ & $1/2$ & $2/3$ & $5/6$ & $11/12$ & $1$ \\
\bottomrule
\end{tabular}
\end{center}

$\mu$'s interior nodes are $t_1,\ldots,t_7$ (plus $\mu(t_0,\cdot)=\rho_0$,
$\mu(t_8,\cdot)=\rho_1$); $\psi$ and the collocated $(\rho,m)$ all live at the $8$ cell
centers $\tau_0,\ldots,\tau_7$ above. (Verified against \code{refine\_ends\_time\_grid(6)}
numerically -- matches to machine precision.)

%------------------------------------------------------------------------------
\section{The discrete FP constraint on this grid}
\label{sec:fp}
%------------------------------------------------------------------------------

The continuous constraint is $\partial_t\rho + \partial_x m = \eps\,\partial_x^2\rho$.
Discretely, the residual is evaluated on the constraint grid but \emph{landed} at the
KE grid's locations $(\tau_i,\xi_k)$ (where $\phi$, the projection's dual potential,
also lives):
\begin{equation}
  \label{eq:residual}
  f(\tau_i,\xi_k) =
  \underbrace{\frac{\mu(t_{i+1},\xi_k)-\mu(t_i,\xi_k)}{\Delta t_i}}_{\partial_t\mu}
  \;+\;
  \underbrace{\frac{\psi(\tau_i,x_k)-\psi(\tau_i,x_{k-1})}{\dx}}_{\partial_x\psi}
  \;-\;
  \eps\,\underbrace{D_x^2\bigl[\bar I_t\mu\bigr](\tau_i,\xi_k)}_{\text{diffusion}},
\end{equation}
where $\bar I_t\mu(\tau_i,\xi_k) := \tfrac12[\mu(t_i,\xi_k)+\mu(t_{i+1},\xi_k)]$ is
exactly the $\hat\rho$ map of \eqref{eq:consensus} (\code{interp\_t\_at\_phi}, reused
verbatim), and $D_x^2$ is the standard $3$-point second difference in $x$
(\code{deriv\_x\_at\_phi} $\circ$ \code{deriv\_x\_at\_m}, with the zero-Neumann ghost at
$k=1,n_x$ -- entirely a space operator, untouched by any of this).

\medskip
\noindent\textbf{Every term in \eqref{eq:residual} is local to cell $i$ in time.} The
$\partial_t\mu$ term only ever reads $\mu$ at $i$'s own two bounding edges, divided by
$i$'s own width -- it never sees a neighboring cell's width, so there is nothing to
special-case at a junction. The diffusion term first applies $\bar I_t$
(Lemma~\ref{lem:midpoint}: exact at $\tau_i$ for \emph{any} $\Delta t_i$, independent of
neighboring cells) and then a purely spatial, uniform-$\dx$ operator. The
$\partial_x\psi$ term never touches time at all. \textbf{Consequence:} the discrete FP
constraint~\eqref{eq:residual}, as literally defined by this codebase, has the
\emph{same} local truncation order at every cell -- including both refined end-cells and
their neighbors -- as it does on the uniform grid. Nothing degrades at the junction; the
refinement changes the (local) resolution, not the (local) accuracy.

\subsection{Worked instantiation, $n_t=6$}

Using the edges from Section~\ref{sec:worked}, at the first junction ($i=0,1,2$):
\begin{align}
  f(\tau_0,\xi_k) &= \frac{\mu(\tfrac1{12},\xi_k)-\rho_0(\xi_k)}{1/12}
    + \frac{\psi(\tau_0,x_k)-\psi(\tau_0,x_{k-1})}{\dx}
    - \eps\, D_x^2\Bigl[\tfrac12\bigl(\rho_0+\mu(\tfrac1{12},\cdot)\bigr)\Bigr](\xi_k), \\
  f(\tau_1,\xi_k) &= \frac{\mu(\tfrac16,\xi_k)-\mu(\tfrac1{12},\xi_k)}{1/12}
    + (\cdots)/\dx
    - \eps\, D_x^2\Bigl[\tfrac12\bigl(\mu(\tfrac1{12},\cdot)+\mu(\tfrac16,\cdot)\bigr)\Bigr](\xi_k), \\
  f(\tau_2,\xi_k) &= \frac{\mu(\tfrac13,\xi_k)-\mu(\tfrac16,\xi_k)}{1/6}
    + (\cdots)/\dx
    - \eps\, D_x^2\Bigl[\tfrac12\bigl(\mu(\tfrac16,\cdot)+\mu(\tfrac13,\cdot)\bigr)\Bigr](\xi_k),
\end{align}
and mirrored at the last junction ($i=5,6,7$):
\begin{align}
  f(\tau_5,\xi_k) &= \frac{\mu(\tfrac56,\xi_k)-\mu(\tfrac23,\xi_k)}{1/6} + (\cdots), \\
  f(\tau_6,\xi_k) &= \frac{\mu(\tfrac{11}{12},\xi_k)-\mu(\tfrac56,\xi_k)}{1/12} + (\cdots), \\
  f(\tau_7,\xi_k) &= \frac{\rho_1(\xi_k)-\mu(\tfrac{11}{12},\xi_k)}{1/12} + (\cdots).
\end{align}
Only the denominator of the $\partial_t\mu$ term changes across each row (per the local
$\Delta t_i$); every other piece of every row has exactly the same form it would on the
uniform grid.

%------------------------------------------------------------------------------
\section{The projection: exact enforcement via the Euclidean adjoint}
\label{sec:projection}
%------------------------------------------------------------------------------

The projection solves $f=0$ \emph{exactly} (up to linear-solve tolerance) by writing the
correction in terms of a potential $\phi$ living at the same $(\tau_i,\xi_k)$ locations
as $f$, and solving $R R^{*}\phi = f$ where $R$ is the linear residual map
\eqref{eq:residual} and $R^{*}$ is \emph{its Euclidean (array-index) adjoint} -- not a
physically/quadrature-weighted adjoint. This is a structural, algebraic fact, not an
approximation: correcting $(\mu,\psi)\mapsto(\mu,\psi) + R^{*}\phi$ gives
$R(\mu+R^*\phi\text{-part},\,\psi+R^*\phi\text{-part}) = f - RR^*\phi = 0$ \emph{exactly},
for \emph{any} $\phi$ solving $RR^*\phi=f$, regardless of whether the individual pieces
of $R^*$ are themselves accurate interpolants of anything physical. This is why grid
non-uniformity can never break \emph{exactness} of the constraint enforcement: exactness
only needs $R^*$ to literally be $R$'s Euclidean transpose, which is a property that is
checked once and for all (\code{scripts/validate\_nu.py}'s adjoint-identity check,
Section~\ref{sec:decisions}), not re-derived per grid.

\subsection{\code{deriv\_t\_at\_rho}: the one place the width jump appears in a formula}

$R^{*}$'s time-derivative piece, $\partial_t^{*}$, is the adjoint of $\partial_t\mu$'s
$1/\Delta t_i$-weighted first difference:
\begin{equation}
  (\partial_t^{*}v)_j = \frac{v_{j+1}}{\Delta t_{j+1}} - \frac{v_j}{\Delta t_j}.
\end{equation}
At the fine/coarse junction $j=1$ (interior node between cell $1$, width $\h/2$, and cell
$2$, width $\h$, in the worked example):
\begin{equation}
  (\partial_t^{*}v)_1 = \frac{v_2}{\h} - \frac{2v_1}{\h},
\end{equation}
a factor-of-$2$ asymmetry between the two terms coming purely from cell $1$ being half
the width of cell $2$ -- each term rescaled by its \emph{own} cell's width, not by any
shared center-to-center distance. Mirrored at $j=5$: $(\partial_t^{*}v)_5 = 2v_6/\h -
v_5/\h$.

\subsection{\code{interp\_t\_at\_rho}: an accuracy caveat (not a correctness one)}

The other time-direction piece used inside the correction, \code{interp\_t\_at\_rho}
(feeding $\nabla_x\phi$'s contribution back into $\rho_{\mathrm{out}}$), is the plain
average $\tfrac12(v_i+v_{i+1})$ mapping cell-center data back to the edge $t_{i+1}$
between them. Unlike Lemma~\ref{lem:midpoint} (which places the average \emph{at} the
midpoint of the two sample points, exactly), this map estimates the value \emph{at}
$t_{i+1}$ from samples at $\tau_i,\tau_{i+1}$ -- and $t_{i+1}$ is the midpoint of
$[\tau_i,\tau_{i+1}]$ only when $\Delta t_i=\Delta t_{i+1}$ (locally uniform). A Taylor
expansion gives
\begin{equation}
  \tfrac12\bigl[g(\tau_i)+g(\tau_{i+1})\bigr] - g(t_{i+1})
  = \frac{\Delta t_{i+1}-\Delta t_i}{4}\,g'(t_{i+1}) + O(\dt^2),
\end{equation}
i.e.\ a first-order error proportional to the local width \emph{mismatch}, vanishing
whenever neighboring cells share a width. At the $n_t=6$ junction ($j=1$: $\Delta
t_1=\h/2\to\Delta t_2=\h$) this is $O(\h)$, localized to that one edge's estimate. This
does \emph{not} affect whether the constraint $f=0$ is satisfied after projection
(Section~\ref{sec:projection} above -- that's exact, algebraically) or the accuracy of
the constraint itself (Section~\ref{sec:fp} -- \code{interp\_t\_at\_rho} never appears
in $f$'s definition, only in the least-squares correction machinery); it is purely a
statement that the \emph{internal} quantity $\nabla_x\phi$ used to build that
correction is a slightly less accurate estimate right at the two junction edges. Since
$\phi$ itself carries no independent physical meaning (it is a Lagrange multiplier, not
part of the state), this has no observable consequence beyond a highly localized
constant in the projection's own internal bookkeeping.

%------------------------------------------------------------------------------
\section{What does not need re-deriving}
%------------------------------------------------------------------------------

\begin{itemize}
  \item \textbf{\code{projection\_nu.py}}'s per-mode tridiagonal operator ($M_0, M_1,
        M_2$) is assembled \emph{numerically}, by probing \code{deriv\_t\_at\_rho} /
        \code{interp\_t\_at\_rho} / etc.\ with unit vectors, not from a closed-form
        formula. It automatically picks up whatever $\Delta t_i$ the grid has, including
        the jumps above -- no new algebra needed.
  \item \textbf{\code{prox\_nu.py}}'s $\sigma_{\mathrm{eff}} = \sigma\, \Delta t_i /
        \dt_{\mathrm{ref}}$ weighting is already generic and diagonal in $i$ (the KE
        prox is pointwise in $(\tau_i,\xi_k)$, so it never sees a neighboring row's
        width at all). In the $n_t=6$ example, $\dt_{\mathrm{ref}} = 1/8$ (mean width of
        the actual $8$-cell grid), so the fine end cells get $\sigma_{\mathrm{eff}} =
        \sigma\cdot\tfrac{1/12}{1/8} = \tfrac23\sigma$ and the coarse middle cells get
        $\sigma_{\mathrm{eff}} = \sigma\cdot\tfrac{1/6}{1/8} = \tfrac43\sigma$.
  \item \textbf{\code{pipeline\_nu.py}}'s $\Delta t^{\mathrm{dual}}_j =
        \tfrac12(\Delta t_j + \Delta t_{j+1})$ quadrature weight for the staggered-grid
        norm is likewise generic; at the junction $\Delta t^{\mathrm{dual}}_1 =
        \tfrac12(\h/2+\h) = \tfrac34\h$.
\end{itemize}

%------------------------------------------------------------------------------
\section{Implementation}
\label{sec:decisions}
%------------------------------------------------------------------------------

Added to \code{time\_grid.py}:
\begin{lstlisting}[language=Python]
def refine_ends_time_grid(nt: int) -> TimeGrid:
    """Uniform nt-cell grid with cell 0 and cell nt-1 each split into two
    equal halves. Returns a TimeGrid with nt+2 cells total."""
    if nt < 2:
        raise ValueError("nt must be >= 2 to refine both ends without overlap")
    h = 1.0 / nt
    base_edges = np.linspace(0.0, 1.0, nt + 1)
    t_edges = np.concatenate([
        [base_edges[0], base_edges[0] + 0.5 * h],
        base_edges[1:-1],
        [base_edges[-1] - 0.5 * h, base_edges[-1]],
    ])
    return _time_grid_from_edges(t_edges)
\end{lstlisting}
Verified against Table 1 by hand: \code{refine\_ends\_time\_grid(6)} reproduces the
$t$/$\Delta t$/$\tau$ values above exactly.

\subsection*{Decisions made on the open design questions}
\begin{enumerate}
  \item \textbf{$n_t$ is the base resolution}, so the returned grid has $n_t+2$ cells,
        not $n_t$. This breaks the ``constructor returns exactly $n_t$ cells'' pattern
        that \code{clustered\_time\_grid} has, but composes more cleanly ($n_t=6$
        unambiguously means ``start from the $6$-cell uniform grid and refine its
        ends''). $n_t < 2$ raises \code{ValueError} (the two end cells would coincide).
  \item \textbf{Fixed single halving, no \code{levels}/\code{k} parameter yet} -- kept to
        exactly what was asked (``start simple''). The name leaves room to add that
        later without an API break.
  \item \textbf{Validation added}: \code{scripts/validate\_nu.py}'s adjoint-identity
        check and repeated-projection stability check both now also run against
        \code{refine\_ends\_time\_grid(NT)}, alongside the existing uniform/clustered
        cases. All checks pass ($5/5$). Not yet added: a dedicated
        \code{sb\_gaussian\_nonuniform.py}-style run against the closed-form SB solution
        on this specific grid.
\end{enumerate}

\end{document}
